\documentclass[10pt,a4paper]{amsart}
\usepackage[margin=19mm]{geometry}
\usepackage{amsmath,amssymb,mathtools}
\usepackage[scr=rsfs,cal=euler]{mathalfa}
\usepackage{xfrac}
\usepackage[unicode,hidelinks,bookmarks=false]{hyperref}
\newcommand{\ie}{i.e.\ }
\newcommand{\cf}{cf.\ }
\newcommand{\resp}{respectively\ }
\newcommand{\Subsection}{\subsection}
\providecommand{\nobreakdash}{\nobreak\hbox{-}}
\setlength{\emergencystretch}{2em}
\multlinegap0pt
\usepackage{amssymb}
\usepackage[shortlabels]{enumitem}
\newtheorem{theorem}{Theorem}
\title{On the maximum number of tangencies among $1$-intersecting curves}
\author{Eyal Ackerman, Bal\'azs Keszegh}
\date{Source: arXiv:2603.11885v1 (CC BY 4.0)}
\begin{document}
\maketitle

\noindent\emph{Source and licence.} On the maximum number of tangencies among $1$-intersecting curves. Version arXiv:2603.11885v1 (CC BY 4.0), distributed by arXiv under the Creative Commons Attribution 4.0 International licence, http://creativecommons.org/licenses/by/4.0/, as recorded in the arXiv licence metadata for this exact version and shown on the abstract page https://arxiv.org/abs/2603.11885v1. Full licence notice: \texttt{licenses/arxiv-2603.11885-CC-BY-4.0.txt}.\par\smallskip

\noindent\textit{Editorial scope.} Counted unit: the article's theorem thm:f-sparse-neighborhoods (source0.tex lines 937--942). The article's complete original proof of it is delivered with it (source0.tex lines 944--986, one \texttt{proof} environment of 43 lines). No mathematics is written, repaired or re-derived: the statement and the proof are the article's own lines, in the article's own notation, and no retained line is substituted or re-flowed.  No further source-local context is needed: every cross-reference of every retained line resolves inside the two delivered spans.  Nothing was omitted from any retained span, so the package carries no omission record.  No retained line contains a citation, so the wrapper carries no bibliography and no bibliography entry.  No retained line contains a figure environment, an image-inclusion command or drawing code, so no image is delivered and the package declares no asset dependency.  Editorial formatting only, in addition: an AMS \texttt{amsart} preamble that restates the article's own declarations and macros used by the retained lines, namely the environments \texttt{theorem} on separate counters, together with the standard packages those lines need.  The retained environments print in this document as Theorem 1 in order of appearance, whereas the article prints the same items with the numbering of its own full document.  The complete native source of the article as distributed in the arXiv e-print for this exact version is bundled unchanged as \texttt{sources/196241/source0.tex}.
\begin{theorem}
	\label{thm:f-sparse-neighborhoods}
	Let $f:\mathbb{R} \rightarrow \mathbb{R}$ be a nonnegative and nondecreasing function and let $G=(A \cup B,E)$ be an $n$-vertex bipartite graph with $\overline{d}(G) \ge 16$. \\
	(i)~If every pair of vertices in $G$ has $f$-sparse sub-bineighborhoods, then $(\overline{d}(G))^3 \le 2^{11}nf(2\overline{d}(G))$.\\
	(ii)~If every pair of adjacent vertices has $f$-sparse sub-bineighborhoods and $|E| \ge 25n^{3/2}$, then $(\overline{d}(G))^3 \le 2^{20}nf(2\overline{d}(G))$.
\end{theorem}

\begin{proof}
	Set $m=|E|$ and let $d=\overline{d}(G)=2m/n$ be the average degree of $G$.
	We first replace $G$ with a graph $G'$ in which the degree of every vertex is roughly the average degree (while maintaining the $f$-sparseness property):
	Go over the vertices of $G$ in an arbitrary order and replace each vertex $v$ such that $\deg(v)>d$ by $\lfloor \deg(v)/d \rfloor$ vertices of degree $d$ and possibly one vertex of a smaller degree.
	The neighbors of $v$ are connected to the vertices that replace $v$ in an arbitrary manner which satisfies this property.
	Observe that the resulting graph still has $f$-sparse sub-bineighborhoods for every pair of (adjacent) vertices, since each sub-neighborhood of a vertex in the new graph cannot contain two vertices that correspond to the same original vertex.
	Thus, every sub-neighborhood in the new graph corresponds to a sub-neighborhood in $G$.
	%Indeed, let $H'$ be such a subgraph.
	%Consider the corresponding subgraph of $G$, denote it by $H$, where each vertex $x'\in V(H')$ is replaced by its corresponding vertex $x \in A \cup B$ (that is, $x$ was split into several vertices one of which is $x'$) and $E(H) = \{ (x,y) \mid (x',y') \in E(H')\}$.
	%Then $H'$ is isomorphic to $H$ and therefore has $f$-sparse neighborhoods.		
	Furthermore, since $d=2m/n$, the number of vertices in the new graph is at most $\sum_{v \in A\cup B} (\lfloor \deg_{G}(v)/d \rfloor + 1) \le 2n$. 
	Because the number of edges remains the same, the average degree of this graph is at least $d/2$.
	
	Next, we repeatedly remove vertices of degree smaller than $d/8$ until no such vertices remain.
	We claim that at least $m/2$ edges remain after this step.
	Indeed, suppose for contradiction that less than $m/2$ edges remain and denote by $x$ the number of vertices that were removed.
	Then at most $xd/8$ edges were removed, therefore $xd/8>m/2=2nd/8$ which is impossible since $x \le 2n$.	
	Denote by $G'=(A' \cup B', E')$ the resulting graph and note that $|A' \cup B'| \le 2n$, $|E'| \ge m/2$ and the degree of every vertex in $G'$ is at least $d/8$ and at most $d$.
	Observe also that  $|A' \cup B'| \ge 2|E'|/d \ge m/d = n/2$.
	Furthermore, $G'$ still has the $f$-sparse sub-bineighborhoods property, since this property is maintained when deleting a vertex.
	
	Considering (i), for each pair of vertices $a \in A'$ and $b \in B'$ the number of edges in the subgraph of $G'$ induced by the union of their open bineighborhoods is at most $f(2d)$.
	Therefore we have $\sum_{a \in A', b \in B'}  |E(G'[(N_{G'}(a)\setminus \{b\}) \cup (N_{G'}(b) \setminus \{a\}])| \le \binom{2n}{2} f(2d) \le 2n^2 f(2d)$.
	Note that every edge $(a,b) \in E'$ is counted in this sum for every pair $a' \in A'$ and $b' \in B'$ such that $a \in N_{G'}(b') \setminus \{a'\}$ and $b \in N_{G'}(a') \setminus \{b'\}$. Since $a$ (resp., $b$) has at least $\frac{d}{8}-1 \ge \frac{d}{16}$ neighbors different from $b$ (resp., $a$), it follows that $(a,b)$ is counted at least $d^2/256$ times in the above sum (here we used the assumption $d \ge 16$).
	Therefore, $nd/2=m=|E| \le 2|E'| \le 2^{10} n^2f(2d)/d^2$ and thus $d^3 \le 2^{11}nf(2d)$.
	
	Considering (ii), we have $d\ge 50n^{1/2}$ since $m\ge 25n^{3/2}$.  Assume w.l.o.g.\ that $|B'|\ge |A'|$ 
	and therefore $|B'| \ge n/4$ and $|A'| \le n$.
	Let $\#K_{2,1}$ be the number of $K_{2,1}$'s in $G'$ with two vertices in $A'$ and 
	let $\#K_{2,2}$ be the number of $K_{2,2}$'s in $G'$.
	Then,
	$$\#K_{2,2} =  \frac{1}{4}\sum_{a \in A', b \in B', (a,b) \in E'}  |E(G'[(N_{G'}(a)\setminus \{b\})\cup (N_{G'}(b)\setminus\{a\}])| \le \frac{m}{4}f(2d) = \frac{dn}{8} f(2d).$$ 
	
	On the other hand, denoting $s_{a,a'}=|N_{G'}(a)\cap N_{G'}(a')|$ for two distinct vertices $a, a'\in A'$ we have $\#K_{2,1} = \sum_{a\ne a' \in A'} s_{a,a'}$ and:
	$$\#K_{2,2}=\sum_{a\ne a'\in A'}\binom{s_{a,a'}}{2}=\frac{1}{2}\sum_{a\ne a'\in A'}(s_{a,a'}^2-s_{a,a'}) \ge \frac{(\#K_{2,1})^2}{2\binom{|A'|}{2}}-\frac{\#K_{2,1}}{2}\ge \frac{(\#K_{2,1})^2}{2\binom{n}{2}}-\frac{\#K_{2,1}}{2},$$
	where the second to last inequality follows from Cauchy's inequality. 
	For $\#K_{2,1}$ we have the following lower bound:	
	
	$$\#K_{2,1}=\sum_{a\ne a'\in A'}s_{a,a'} = \sum_{b \in B'} \binom{\deg(b)}{2} \ge \frac{n}{4}\binom{d/8}{2}\ge nd^2/2^{11}\ge n^2,$$
	where we used double-counting by the vertex from $B'$ for the $K_{2,1}$'s and the inequalities $|B'|\ge n/4$ and $d\ge 50n^{1/2}$. Combining the inequalities above we get:
	$$\frac{dnf(2d)}{8} \ge \#K_{2,2} \ge \frac{(\#K_{2,1})^2}{n^2}-\frac{\#K_{2,1}}{2} \ge \frac{(\#K_{2,1})^2}{n^2}-\frac{(\#K_{2,1})^2}{2n^2} = \frac{(\#K_{2,1})^2}{2n^2} \ge \frac{n^2d^4}{2^{22}\cdot 2n^2}=\frac{d^4}{2^{23}}.$$	
	This implies that $d^3\le 2^{20}nf(2d)$, as claimed.
\end{proof}

\end{document}
